\documentclass[12pt,a4paper,bibliography=totocnumbered,listof=totocnumbered]{scrartcl}
\usepackage[ngerman]{babel}
\usepackage[utf8]{inputenc}
\usepackage{amsmath}
\usepackage{amsfonts}
\usepackage{amssymb}
\usepackage{graphicx}
\usepackage{fancyhdr}
\usepackage{tabularx}
\usepackage{geometry}
\usepackage{setspace}
\usepackage[right]{eurosym}
\usepackage[printonlyused]{acronym}
\usepackage{subfig}
\usepackage{floatflt}
\usepackage[usenames,dvipsnames]{color}
\usepackage{colortbl}
\usepackage{paralist}
\usepackage{array}
\usepackage{titlesec}
\usepackage{parskip}
\usepackage[right]{eurosym}
\usepackage{picins}
\usepackage[subfigure,titles]{tocloft}
\usepackage[pdfpagelabels=true]{hyperref}

\usepackage[
	backend=biber,
	style=alphabetic,
	citestyle=alphabetic,
	sorting=nty,
	doi=true,url=true
]{biblatex}
\bibliography{bib}

\usepackage{listings}
\lstset{basicstyle=\footnotesize, captionpos=b, breaklines=true, showstringspaces=false, tabsize=2, frame=lines, numbers=left, numberstyle=\tiny, xleftmargin=2em, framexleftmargin=2em}
\makeatletter
\def\l@lstlisting#1#2{\@dottedtocline{1}{0em}{1em}{\hspace{1,5em} Lst. #1}{#2}}
\makeatother

\geometry{a4paper, top=27mm, left=30mm, right=20mm, bottom=35mm, headsep=10mm, footskip=12mm}

\hypersetup{unicode=false, pdftoolbar=true, pdfmenubar=true, pdffitwindow=false, pdfstartview={FitH},
	pdftitle={Abschlussarbeit},
	pdfauthor={Daniel Brettschneider},
	pdfsubject={Abschlussarbeit},
	pdfcreator={\LaTeX\ with package \flqq hyperref\frqq},
	pdfproducer={pdfTeX \the\pdftexversion.\pdftexrevision},
	pdfkeywords={Abschlussarbeit},
	pdfnewwindow=true,
	colorlinks=true,linkcolor=black,citecolor=black,filecolor=magenta,urlcolor=black}
\pdfinfo{/CreationDate (D:20110620133321)}

% Benutzerdefinierte Befehle für späteren Werteabruf
\def\title#1{\gdef\@title{#1}\gdef\THETITLE{#1}}
\def\author#1{\gdef\@author{#1}\gdef\THEAUTHOR{#1}}
\newcommand{\authoremail}{\href{mailto:\email}{\email}}

% Wertezuweisung in Variablen, damit diese nur an einer Stelle geändert werden müssen
\title{Entwicklung eines Android Usability Testing Frameworks}
\author{Albert Hoffmann}
\newcommand{\email}{albert.hoffmann@hs-osnabrueck.de}
\newcommand{\pruefer}{Prof. Michaela Ramm, M.A.}

\begin{document}

\titlespacing{\section}{0pt}{12pt plus 4pt minus 2pt}{-6pt plus 2pt minus 2pt}

% Kopf- und Fusszeile
\renewcommand{\sectionmark}[1]{\markright{#1}}
\renewcommand{\leftmark}{\rightmark}
\pagestyle{fancy}
\lhead{}
\chead{}
\rhead{\thesection\space\contentsname}
\lfoot{\THETITLE}
\cfoot{}
\rfoot{\thepage}
\renewcommand{\headrulewidth}{0.4pt}
\renewcommand{\footrulewidth}{0.4pt}

% Vorspann
\renewcommand{\thesection}{\Roman{section}}
\renewcommand{\theHsection}{\Roman{section}}
\pagenumbering{Roman}

% ----------------------------------------------------------------------------------------------------------
% Titelseite
% ----------------------------------------------------------------------------------------------------------


\thispagestyle{empty}
\begin{center}
	\includegraphics[scale=1]{img/hs_os}\\
	\vspace*{2cm}
	\Large
	\textbf{Fakultät}\\
	\textbf{Ingenieurwissenschaften und Informatik}\\
	\vspace*{2cm}
	\Huge
	\textbf{Projektdokumentation}\\
	\vspace*{0.5cm}
	\large
	für das \\
	\textbf{Masterprojekt}\\  
	mit dem Thema\\
	\vspace*{1cm}
	\textbf{\THETITLE}\\
	\vspace*{2cm}
	
	\vfill
	\normalsize
	\newcolumntype{x}[1]{>{\raggedleft\arraybackslash\hspace{0pt}}p{#1}}
	\begin{tabular}{x{6cm}p{7.5cm}}
		\rule{0mm}{5ex}\textbf{Autor:} & \THEAUTHOR \newline 
		\authoremail \\ 
		\rule{0mm}{5ex}\textbf{Prüfer:} & \pruefer \\ 
		\rule{0mm}{5ex}\textbf{Abgabedatum:} & 06.10.2014 \\ 
	\end{tabular} 
\end{center}
\pagebreak


\begin{abstract}

\noindent\textbf{Zusammenfassung:}\\ \\
Die vorliegende Ausarbeitung wurde in \LaTeX verfasst und ist eine gemeinsame Arbeit von Albert Hoffmann und Oliver Erxleben an der Hochschule Osnabrück / University of Applied Sciences im Masterstudiengang Informatik - Mobile und Verteilte Anwendungen des Fachbereichs Ingenieurswissenschaften und Informatik für das Modul Mensch-Maschine-Kommunikation im Wintersemester 2013/14. Die Arbeit dokumentiert die Entwicklung einer Programmierbibliothek zum Test der Usability von Mobilanwendungen für das Android-Betriebssystem.\\ \\
Die Arbeit gliedert sich in verschiedene Teile.\\
%TODO: update Aufbau und Ablauf
\end{abstract}
% Verzeichnisse
% ----------------------------------------------------------------------------------------------------------
% Verzeichnisse
% ----------------------------------------------------------------------------------------------------------
% TODO Typ vor Nummer
\renewcommand{\cfttabpresnum}{Tab. }
\renewcommand{\cftfigpresnum}{Abb. }
\settowidth{\cfttabnumwidth}{Abb. 10\quad}
\settowidth{\cftfignumwidth}{Abb. 10\quad}

\titlespacing{\section}{0pt}{12pt plus 4pt minus 2pt}{2pt plus 2pt minus 2pt}
\singlespacing
\rhead{INHALTSVERZEICHNIS}
\renewcommand{\contentsname}{II Inhaltsverzeichnis}
\phantomsection
\addcontentsline{toc}{section}{\texorpdfstring{II \hspace{0.35em}Inhaltsverzeichnis}{Inhaltsverzeichnis}}
\addtocounter{section}{1}
\tableofcontents
\pagebreak
\rhead{VERZEICHNISSE}
\listoffigures
\pagebreak
\listoftables
%\pagebreak
\renewcommand{\lstlistlistingname}{Listing-Verzeichnis}
{\labelsep2cm\lstlistoflistings}
\pagebreak
% ----------------------------------------------------------------------------------------------------------
% Abkürzungen
% ----------------------------------------------------------------------------------------------------------
\section{Abkürzungsverzeichnis}
%TODO: längste Abkürzung steht in eckigen Klammern
\begin{acronym}[HATEOAS]
\setlength{\itemsep}{-\parsep} % geringerer Zeilenabstand
\acro{OSGi}{Open Service Gateway initiative} %TODO entfernen
\acro{HTTP}{Hypertext Transfer Protocol}
\acro{SSL}{Secure Sockets Layer}
\acro{TLS}{Transport Layer Security}
\acro{UI}{User Interface}
\acro{JSON}{JavaScript Object Notation}
\acro{UX}{User Experience}
\acro{SDK}{Software Development Kit}
\acro{IDE}{Integrated Development Environment}
\acro{API}{Application Programming Interface}
\acro{NPM}{Node Package Manager}
\acro{VM}{Virtuelle Maschine}
\acro{SSH}{Secure Shell}
\acro{JAR}{Java Archive File}
\acro{CSS}{Cascading Style Sheets}
\acro{SQL}{Structured Query Language}
\acro{ER}{Entity Relationship}
\acro{REST}{Representational State Transfer}
\acro{URL}{Uniform Resource Locator}
\acro{SHA}{Secure Hash Algorithm}
\acro{AES}{Advanced Encryption Standard}
\acro{AS}{Application Server}
\acro{JPA}{Java Persistence API}
\acro{DBMS}{Database Management System}
\acro{JDBC}{Java Database Connectivity}
\acro{JAX-RS}{Java API for RESTful Web Services}
\acro{RPC}{Remote Procedure Call}
\acro{HATEOAS}{Hypermedia as the Engine of Application State}
\acro{CRUD}{Create, read, update and delete}
\acro{DAO}{Data Access Object}
\acroplural{DAO}[DAO's]{Data Access Objects}
\end{acronym}


% ----------------------------------------------------------------------------------------------------------
% Inhalt
% ----------------------------------------------------------------------------------------------------------
% Abstände Überschrift
\titlespacing{\section}{0pt}{12pt plus 4pt minus 2pt}{-6pt plus 2pt minus 2pt}
\titlespacing{\subsection}{0pt}{12pt plus 4pt minus 2pt}{-6pt plus 2pt minus 2pt}
\titlespacing{\subsubsection}{0pt}{12pt plus 4pt minus 2pt}{-6pt plus 2pt minus 2pt}

% Kopfzeile
\renewcommand{\sectionmark}[1]{\markright{#1}}
\renewcommand{\subsectionmark}[1]{}
\renewcommand{\subsubsectionmark}[1]{}
\lhead{Kapitel \thesection}
\rhead{\rightmark}

\onehalfspacing
\renewcommand{\thesection}{\arabic{section}}
\renewcommand{\theHsection}{\arabic{section}}
\setcounter{section}{0}
\pagenumbering{arabic}
\setcounter{page}{1}


\section{Einleitung}
\subsection{Bezug der Programmquellen und Dokumentation}
Die Programmquellen der entwickelten, bzw. weiterentwickelten Software ist in dem folgenden \emph{GitHub}-Repository abrufbar:
\begin{center}
	\url{https://github.com/Chromdioxyde/hipsterbility}
\end{center}
Als Versionsverwaltung wird \emph{Git}\footnote{Git Webseite: \url{http://git-scm.com/}} eingesetzt.
Mit dem Befehl
\\
 \texttt{git clone https://github.com/Chromdioxyde/hipsterbility.git} kann eine lokale Kopie erstellt werden.
Alternativ kann der Inhalt des Repositories auf der o.~g. Webseite auch als ZIP-Datei heruntergeladen werden.
Die Quellen der Programme liegen im \emph{IntelliJ IDEA}\footnote{IntelliJ IDEA Website: \url{https://www.jetbrains.com/idea/}} Projektformat vor.
Für die Entwicklung wurde \emph{IntelliJ IDEA Ultimate 13.1.5} verwendet.

\subsection{Motivation}
Bei begrenztem Budget, z.~B. von Studierenden sind umfangreiche Labortests oder das Beauftragen von Dienstleistern oft nicht möglich.

Ein weiterer Faktor ist auch die Zielgruppe für die entwickelte Anwendung.
Bei räumlich verteilten Anwendern sind Labortests durch den logistischen Aufwand nur schwer möglich und auch die finanzielle Belastung steigt durch Transportkosten und evtl. mobiles Equipment.

Deshalb wurde im Wintersemester 2013/2014, zusammen mit Oliver Erxleben, ein Projekt durchgeführt, mit dem Ziel Usability-Tests mit minimalem Aufwand für den Teilnehmer durchzuführen (siehe Abschnitt \ref{subsec:stand_bei_projektbeginn}).

Die damals gesammelten Erkenntnisse und Erfahrungen sollen nun vertieft werden, mit dem Ziel die bestehenden Prototypen zu einem nutzbaren Framework weiter zu entwickeln.


\subsection{Zielsetzung}
Das Ziel ist es, ein Usability-Testing Framework oder Toolkit zu entwickeln, mit dem auch einzelne Entwickler oder kleine Teams ihre Anwendungen und Prototypen mit der Hilfe von Testbenutzern testen können.
Für den eigentlichen Testablauf beim Benutzers soll nur ein Android basiertes Gerät und eine Internetverbindung benötigt werden.
Ziel dabei ist es, das Testen möglichst einfach zu gestalten, ohne Labor und zusätzliches Equipment.

Auf der Seite des Testleiters soll nur ein PC oder Notebook benötigt werden, ohne zusätzliche Hardware.
Die Serverkomponente wurde ausgelagert, für den Fall, dass ein Server zur Verfügung steht, bzw. z.~B. von einer Bildungseinrichtung bereitgestellt wird.
\subsection{Problemstellung}
%TODO: Problemstellung
\subsection{Stand der Technik}

\subsubsection{Mobile Usability Testing Methoden}

\subsubsection{Mobile Usability Testing Anbieter}

\subsubsection{Mobile Usability Testing Werkzeuge}



\section{Konzept}

\section{Realisierung}

\section{Ergebnisse}

\section{Fazit}

% ----------------------------------------------------------------------------------------------------------
% Literatur
% ----------------------------------------------------------------------------------------------------------
\renewcommand\refname{Quellenverzeichnis}
\pagebreak

% ----------------------------------------------------------------------------------------------------------
% Anhang
% ----------------------------------------------------------------------------------------------------------
\pagenumbering{Roman}
\setcounter{page}{1}
\lhead{Anhang \thesection}

\begin{appendix}
\section*{Anhang}
\phantomsection
\addcontentsline{toc}{section}{Anhang}
\addtocontents{toc}{\vspace{-0.5em}}

\section{GUI}
Ein toller Anhang.

\subsection*{Screenshot}
\label{app:screenshot}
Unterkategorie, die nicht im Inhaltsverzeichnis auftaucht.

\end{appendix}


\newpage
\thispagestyle{empty}
\begin{center}
	\vspace*{5em}
	\huge\textbf{Erklärung}\\
\end{center}
\vspace{2em}
Hiermit versichere ich, dass ich meine Abschlussarbeit selbständig verfasst und keine anderen als die angegebenen Quellen und Hilfsmittel benutzt habe.

\vspace{4em}
\begin{minipage}{\linewidth}
	\begin{tabular}{p{15em}p{15em}}
		Datum: &  .......................................................\\
		& \centering (Unterschrift)\\
	\end{tabular}
\end{minipage}

\end{document}
